\clearpage
\VerdierPhysicalPageAnchor{0185}%
\begingroup\small\setlength{\parskip}{0pt}\setlength{\emergencystretch}{1.5em}
\setlength{\abovedisplayskip}{4pt plus 2pt minus 2pt}%
\setlength{\belowdisplayskip}{4pt plus 2pt minus 2pt}%

\VerdierSemanticAnchor{ch3d-1-3-6-conditions-cont}
\noindent\textit{i) $F$ factors through $K(\mathcal A)$.}

\noindent\textit{ii) $F$ sends homotopy equivalences to isomorphisms in
$C$.}

\VerdierSemanticAnchor{ch3d-1-3-6-proof}
It is clear that i) $\Rightarrow$ ii). To prove that ii)
$\Rightarrow$ i), it suffices to show that $F$ sends pairs of homotopic
morphisms to equal morphisms. Let $X$ be an object of
$\operatorname{comp}(\mathcal A)$. The functor on
$\operatorname{comp}(\mathcal A)$ given by
\[
Y\longmapsto
\{\text{pairs of morphisms from $X$ to $Y$ equipped with a homotopy joining
them}\}
\]
is represented by a complex $Z$:

\VerdierSemanticAnchor{ch3d-1-3-6-representing-complex}
\begin{equation}
\begin{aligned}
Z^i&=X^i\oplus X^i\oplus X^{i+1},\\[2pt]
d_Z^i&=
\begin{pmatrix}
d_X^i&0&\operatorname{id}_{X^{i+1}}\\
0&d_X^i&-\operatorname{id}_{X^{i+1}}\\
0&0&-d_X^{i+1}
\end{pmatrix}.
\end{aligned}
\tag{1.3.6.1}
\end{equation}
There are therefore two homotopic morphisms
$p_1,p_2\colon X\longrightarrow Z$:

\VerdierSemanticAnchor{ch3d-1-3-6-morphisms}
\begin{equation}
p_1^i=
\begin{pmatrix}
\operatorname{id}_{X^i}\\0\\0
\end{pmatrix},
\qquad
p_2^i=
\begin{pmatrix}
0\\\operatorname{id}_{X^i}\\0
\end{pmatrix},
\tag{1.3.6.2}
\end{equation}
and it suffices to prove that $F(p_1)=F(p_2)$. Now the pair of morphisms
$(\operatorname{id}_X,\operatorname{id}_X)$ equipped with the trivial
homotopy defines a morphism $\pi\colon Z\longrightarrow X$:

\VerdierSemanticAnchor{ch3d-1-3-6-retraction}
\begin{equation}
\pi^i=(\operatorname{id}_{X^i},\operatorname{id}_{X^i},0)
\tag{1.3.6.3}
\end{equation}
One has $\pi p_1=\pi p_2=\operatorname{id}_X$. Moreover, one checks
immediately that the morphisms $\pi$, $p_1$, and $p_2$ are homotopy
equivalences. It follows that $F(p_1)$, $F(p_2)$, and $F(\pi)$ are
isomorphisms and that
$F(\pi)F(p_1)=F(\pi)F(p_2)=\operatorname{id}_{F(X)}$. Consequently,
$F(p_1)=F(p_2)$, which completes the proof of the lemma.

\VerdierSemanticAnchor{ch3d-1-3-7}
\medskip\noindent\textbf{Remark 1.3.7.}\quad
It follows from Proposition (1.3.5) that the derived category of
$\mathcal A$ is the solution to the universal problem of inverting the
quasi-isomorphisms in $\operatorname{comp}(\mathcal A)$ (Chap. II, 2.2.3,
b)).
\endgroup
